\documentclass[12pt, twoside]{article}
%\documentclass[12pt, twoside]{article}
\usepackage[letterpaper, margin=1in, headsep=0.2in]{geometry}
\setlength{\headheight}{0.6in}
%\usepackage[english]{babel}
\usepackage[utf8]{inputenc}
\usepackage{microtype}
\usepackage{amsmath}
\usepackage{amssymb}
%\usepackage{amsfonts}
\usepackage[nomessages]{fp} %\FPeval{\var-name}{2*sin(pi/6)}
\usepackage{siunitx} %units in math. eg 20\milli\meter
\usepackage{yhmath} % for arcs, overparenth command
\usepackage{tikz} %graphics
\usetikzlibrary{quotes, angles, arrows, arrows.meta}
\usepackage{pgfplots}
\usepgfplotslibrary{fillbetween}
\pgfplotsset{compat=1.16}
\usepackage{graphicx} %consider setting \graphicspath{{images/}}
\usepackage{parskip} %no paragraph indent
\usepackage{enumitem}
\usepackage{multicol}
\usepackage{venndiagram}

\usepackage{fancyhdr}
\pagestyle{fancy}
\fancyhf{}
\renewcommand{\headrulewidth}{0pt} % disable the underline of the header
\raggedbottom
\hfuzz=2mm %suppresses overfull box warnings

\usepackage{hyperref}
\usepackage{eurosym}

\fancyhead[LE]{\thepage}
\fancyhead[RO]{Markscheme}
\fancyhead[LO]{La Scuola d'Italia / Dr.\ Huson / IB Math 11 \\* 14 October 2026}

\begin{document}
\subsubsection*{1.12 Homework: Review of Sequences and Series --- Markscheme}

\begin{enumerate}[itemsep=0.3cm]

%--- Problem 1: Arithmetic sequence; d, u_n, u_25, which term is 143 [7 marks] ---
\item[1.] [Maximum mark: 7]

$7, 11, 15, 19, \ldots$, so $u_{1} = 7$.

\textbf{(a)} $d = 4$ \hfill \textbf{A1}

\textbf{(b)} $u_{n} = 7 + (n - 1)4$ \hfill \textbf{M1}

$u_{n} = 4n + 3$ \hfill \textbf{A1}

\emph{Accept the unsimplified form $7 + 4(n - 1)$ for both marks.}

\textbf{(c)} $u_{25} = 4(25) + 3 = 103$ \hfill \textbf{A1}

\textbf{(d)} $4n + 3 = 143$ \hfill \textbf{M1}

$4n = 140$ \hfill \textbf{(A1)}

$n = 35$: the $35^{\text{th}}$ term \hfill \textbf{A1}

\emph{M1 for setting the expression for $u_{n}$ equal to $143$. FT from
(b).}

\medskip

%--- Problem 2: Arithmetic sequence from two terms; u_1, d, u_20, S_20 [7 marks] ---
\item[2.] [Maximum mark: 7]

\textbf{(a)} $u_{1} + 2d = 14$ and $u_{1} + 7d = 39$ \hfill \textbf{M1}

$5d = 25$, so $d = 5$ \hfill \textbf{A1}

$u_{1} = 14 - 2(5) = 4$ \hfill \textbf{A1}

\emph{M1 for both equations, or for $\frac{39 - 14}{8 - 3}$. Accept a GDC
simultaneous-equation solve.}

\textbf{(b)} $u_{20} = 4 + 19(5) = 99$ \hfill \textbf{(M1)A1}

\textbf{(c)} $S_{20} = \dfrac{20}{2}(4 + 99)$ \hfill \textbf{M1}

$= 1030$ \hfill \textbf{A1}

\emph{Accept $S_{20} = \frac{20}{2}\left(2(4) + 19(5)\right)$. FT from (a)
and (b).}

\newpage

%--- Problem 3: Arithmetic series with a given last term; number of terms, sum [4 marks] ---
\item[3.] [Maximum mark: 4]

$u_{1} = 5$, $d = 4$, last term $201$.

\textbf{(a)} $5 + (n - 1)4 = 201$ \hfill \textbf{M1}

$n = 50$ \hfill \textbf{A1}

\emph{The common error is $\frac{201 - 5}{4} = 49$, forgetting to add $1$:
M1A0.}

\textbf{(b)} $S_{50} = \dfrac{50}{2}(5 + 201)$ \hfill \textbf{M1}

$= 5150$ \hfill \textbf{A1}

\emph{FT from the candidate's $n$ in (a).}

\medskip

%--- Problem 4: Theatre seats; u_25, S_25, least number of rows for 600 seats [7 marks] ---
\item[4.] [Maximum mark: 7]

$u_{1} = 20$, $d = 2$, $25$ rows.

\textbf{(a)} $u_{25} = 20 + 24(2) = 68$ seats \hfill \textbf{(M1)A1}

\textbf{(b)} $S_{25} = \dfrac{25}{2}(20 + 68)$ \hfill \textbf{M1}

$= 1100$ seats \hfill \textbf{A1}

\emph{FT from (a).}

\textbf{(c)} $\dfrac{n}{2}\left(2(20) + (n - 1)2\right) \geq 600$, \quad
i.e.\ $n(n + 19) \geq 600$ \hfill \textbf{M1}

$S_{16} = 560$ and $S_{17} = 612$ \hfill \textbf{(A1)}

$17$ rows \hfill \textbf{A1}

\emph{M1 for setting up the sum of the first $n$ rows $\geq 600$ (or $= 600$).
The GDC table or the quadratic $n^{2} + 19n - 600 = 0$, giving
$n = 16.77\ldots$, then rounding \emph{up} to $17$. The answer $16$ (rounding
down) earns M1A1A0: $16$ rows hold only $560$ seats.}

\newpage

%--- Problem 5: Geometric sequence; r, u_n, u_10, first term over 10 000 [7 marks] ---
\item[5.] [Maximum mark: 7]

$3, 6, 12, 24, \ldots$, so $u_{1} = 3$.

\textbf{(a)} $r = 2$ \hfill \textbf{A1}

\textbf{(b)} $u_{n} = 3 \times 2^{n - 1}$ \hfill \textbf{A1}

\emph{$6^{n - 1}$ (multiplying $3 \times 2$ first) earns A0.}

\textbf{(c)} $u_{10} = 3 \times 2^{9} = 1536$ \hfill \textbf{(M1)A1}

\textbf{(d)} $3 \times 2^{n - 1} > 10\,000$ \hfill \textbf{M1}

$u_{12} = 6144$, \quad $u_{13} = 12\,288$ \hfill \textbf{(A1)}

The $13^{\text{th}}$ term, $12\,288$ \hfill \textbf{A1}

\emph{M1 for the inequality, a GDC table, or listing terms. Final A1 needs
both the value and the term number.}

\medskip

%--- Problem 6: Geometric sequence from two terms; r, u_1, u_7, S_6 [6 marks] ---
\item[6.] [Maximum mark: 6]

\textbf{(a)} $u_{1}r = 12$ and $u_{1}r^{4} = 324$ \hfill \textbf{M1}

$r^{3} = \dfrac{324}{12} = 27$, so $r = 3$ \hfill \textbf{A1}

$u_{1} = \dfrac{12}{3} = 4$ \hfill \textbf{A1}

\textbf{(b)} $u_{7} = 4 \times 3^{6} = 2916$ \hfill \textbf{A1}

\textbf{(c)} $S_{6} = \dfrac{4\left(3^{6} - 1\right)}{3 - 1}$
\hfill \textbf{M1}

$= \dfrac{4(728)}{2} = 1456$ \hfill \textbf{A1}

\emph{FT from (a). $u_{6} = 972$ (a term, not a sum) earns no marks.}

\newpage

%--- Problem 7: Training distance growing 10% per week; r, week-8 distance, 8-week total [5 marks] ---
\item[7.] [Maximum mark: 5]

\textbf{(a)} $r = 1.1$ \hfill \textbf{A1}

\emph{$r = 0.1$ or $r = 10$ earns A0.}

\textbf{(b)} $u_{8} = 5 \times 1.1^{7}$ \hfill \textbf{(M1)}

$= 9.7435\ldots = 9.74$ km \hfill \textbf{A1}

\textbf{(c)} $S_{8} = \dfrac{5\left(1.1^{8} - 1\right)}{1.1 - 1}$
\hfill \textbf{(M1)}

$= 57.179\ldots = 57.2$ km \hfill \textbf{A1}

\emph{An exponent of $8$ in (b), giving $10.7$, earns M1A0. FT from (a).}

\medskip

%--- Problem 8: Compound interest, quarterly vs annual; interest earned; years to exceed 8000 [8 marks] ---
\item[8.] [Maximum mark: 8]

\emph{\textbf{GDC:} the Finance Solver may be used with $N = 4n$,
$I(\%) = 4$, $PV = -5000$, $PMT = 0$, $P/Y = C/Y = 4$.}

\textbf{(a)} $FV = 5000\left(1 + \dfrac{4}{100 \times 4}\right)^{4 \times 6}
= 5000(1.01)^{24}$ \hfill \textbf{M1}

$= $ \euro$6348.67$ \hfill \textbf{A1}

\emph{Annual compounding in (a) ($6326.60$) earns M0A0.}

\textbf{(b)} $6348.67 - 5000 = $ \euro$1348.67$ \hfill \textbf{A1}

\emph{FT from (a).}

\textbf{(c)} $5000(1.04)^{6} = $ \euro$6326.60$ \hfill \textbf{(M1)A1}

\textbf{(d)} $5000(1.01)^{4n} > 8000$ \hfill \textbf{M1}

After $11$ years: \euro$7746.59$; after $12$ years: \euro$8061.13$
\hfill \textbf{(A1)}

$12$ years \hfill \textbf{A1}

\emph{M1 for the inequality, a GDC table, or the Finance Solver with
$FV = 8000$ ($N = 47.2$ quarters). $11.8$ years, unrounded, earns M1A1A0;
the question asks for complete years.}

\end{enumerate}
\end{document}
